\subsection{負tangent第一割当ての全正Pfaffian分枝も空}
\label{subsec:h3-order2-22-tangent-plus-Ore}
\begin{lemma}[非零片側curvatureの強化された正常核]
\label{lem:h3-order2-22-curved-graded-zero}
$F\in\mathrm{SL}_2$、$\alpha\ne0$ とし、既知curved SL2包絡
\[
 B_0(F)=k\langle x,y,u,v\rangle/
 ([x,y]-\alpha u^2,\ uU-F(U)u,\ vU-F^{-1}(U)v,\ uv+vu)
\]
を考える。
$++--$ の正常三次核は、$\operatorname{tr}F=1$ の場合を除いて零である。
$\operatorname{tr}F=1$ では、$F$ の固有座標で $kx^3+ky^3$ に等しく、特に負変数 $v$ を含まない。
\end{lemma}
\begin{proof}
$F$ が対角化可能なら命題\ref{prop:h3-order2-22-curved-diagonal}の全係数正常核判定である。
$\alpha\ne0$ なので $F=\pm I$ のWeyl中央化核も零となる。
残る固有値は $\lambda^2-\lambda+1=0$ だけで、正常核は表示した二次元。
非対角型 $F=\pm J$ は $u$ の尺度で $\alpha=-1$ に正規化できる。
この二点のPBW正常行列にはそれぞれ非零な定数 $10\times10$ minorがあり、核は零である。
前の全SL2判定では一次因子までしか使わなかったこの二点を、ここでは正常核零まで強化した。
\end{proof}

\begin{theorem}\label{thm:h3-order2-22-tangent-first-semisimple-empty}
負tangent第一割当て
\[
 [x,y]+u^2=0,\qquad q_U+uv+vu=0,\qquad
 q_U=x^2\ \text{又は}\ xy+yx
\]
の全半単純中山H3候補は、両Pfaffian分枝・全退化係数込みで空である。
$q_U$ が非退化の場合は中山のJordan部分も既に排除されているので、全Jordan型込みで空である。
平方型の非半単純対象や、tangent逆割当てをこの定理に含めない。
\end{theorem}
\begin{proof}
負Pfaffianは前節で全値排除済みなので、正Pfaffianを扱う。
全atlasの係数から
\[
 F=\begin{pmatrix}-z_2&z_0\\-z_5&-z_3\end{pmatrix}\in\mathrm{SL}_2
\]
を作ると、この $F$ は全 $\mathrm{SL}_2$ を走る。
三生成元部分を
\[
 S(F)=k\langle U,u\rangle/([x,y]+u^2,\ uU-F(U)u)
\]
と書けば、包絡は $B=S(F)[v;\sigma,\delta]$ であり、
$\sigma|_U=F^{-1},\sigma(u)=-u$。
$F=\left(\begin{smallmatrix}r&s\\t&w\end{smallmatrix}\right)$ に対する微分は、平方型では
\[
 \delta(x)=-\tfrac t2xu,\quad \delta(y)=rxu+\tfrac t2yu,\quad \delta(u)=-x^2;
\]
非退化型では
\[
 \delta(x)=-\tfrac{r+w}{2}xu-tyu,\quad
 \delta(y)=sxu+\tfrac{r+w}{2}yu,\quad \delta(u)=-(xy+yx).
\]
この表示は $z_5=0$ 等の境界を除去せず成立する。
全 $\det F=1$ で各四つのoverlapが解消することも検証する。
$v$ 次数filtrationのassociated gradedは、補題の $B_0(F)$、$\alpha=-1$ そのものである。

元の非零正常三次 $G$ があれば、その $v$ 最高次数項は $B_0(F)$ の正常核に入る。
補題より $\operatorname{tr}F=1$ に限り、しかも最高 $v$ 次数は0。
従って $G$ 自身が $S(F)$ に属し、固有座標では $G=Ax^3+By^3$ である。
ここで残った正常核を零と取り違えてはいけない。
$AB=0$ なら既に一次因子を持つ。
$AB\ne0$ では定理\ref{thm:h3-order2-22-primitive6-nondomain}の三次norm構成が
三生成元部分 $S(F)/(G)$ 自身に適用され、零因子を与える。
同定理の $\beta=0$ の場合、構成に負変数 $v$ は不要であり、
$u$ と一方の固有座標を逆転したGWAで
$\sigma^3a_0=a_0$ と $X^6=\rho^3\operatorname{Norm}_3(a_0)$ から
$Z=(\rho a_0)^{-1}X^2$、$Z^3=1$ を得る。
monic六基底により $Z-1,Z^2+Z+1$ はともに非零である。
従ってこの三生成元商はdomainでない。

$Gv+vG=0$ をOre表示で比較すると $\sigma(G)=-G,\delta(G)=0$。
従って
\[
 B/(G)=(S(F)/(G))[v;\bar\sigma,\bar\delta]
\]
はmonic $v$ 基底を持ち、$S(F)/(G)$ が単射で入る。
その零因子は消えないため、全 $B/(G)$ も非domainでH3に反する。
これはtwist不変性やfiltered domain性の逆推論を用いない。
最後の全Jordan型の主張は非退化正topのJordan排除定理による。
\end{proof}

全SL2の二組のPBW恒等式と associated graded の二つのJordan正常minorは
\[
\texttt{verify\_cubic\_order2\_22\_tangent\_plus\_Ore.py}
\]
で再生する。二パラメータ一般点のminorだけでは見落とす $\operatorname{tr}F=1$ の二次元正常核を、
この証明では三生成元商の非整域性によって全て排除している。
